% ថ្នាក់ទី ១១ (វិទ្យាសាស្ត្រ) · ជំពូកទី ១ · មេរៀនទី ៥ · សៀវភៅទំព័រ ៣៦–៤៣
\chapter{ទំនាក់ទំនងតួនៃស៊ីត}

\begin{objBox}[វត្ថុបំណង]
\begin{enumerate}[1]
  \item កំណត់តួទី $n$ ដោយប្រើស៊ីតជំនួយ
  \item កំណត់តួទី $n$ តាមទំនាក់ទំនងរវាង $a_n$ និង $S_n$
  \item កំណត់តួទូទៅនៃទំនាក់ទំនងតួក្នុងទម្រង់ $a_{n+2} + pa_{n+1} + qa_n = 0$ ។
\end{enumerate}
\end{objBox}

\section{កំណត់តួទី \texorpdfstring{$n$}{n} ដោយប្រើស៊ីតជំនួយ}

គេបានសិក្សារួចមកហើយនូវទំនាក់ទំនងតួនៃស៊ីតនព្វន្ត និងស៊ីតធរណីមាត្រ ។

ចំពោះស៊ីតនព្វន្ត $a_{n+1} = a_n + d$, $d$ ជាផលសងរួម ។

ចំពោះស៊ីតធរណីមាត្រ $a_{n+1} = a_n \times q$, $q$ ជាផលធៀបរួម ។

តទៅទៀតគេនឹងសិក្សាអំពីទំនាក់ទំនងតួនៃស៊ីតផ្សេងៗទៀតដោយកំណត់រកតួទី $n$ នៃស៊ីតនោះ ។

\begin{example}[លំហាត់គំរូ ១]
កំណត់តួទូទៅនៃស៊ីត $(a_n)$ ដោយទំនាក់ទំនងកំណើន $a_1 = 1$, $a_{n+1} = 1 + 4a_n$ ។
\end{example}
\ans
\textbf{របៀបទី ១} ៖ គេមាន $a_1 = 1$, $a_{n+1} = 1 + 4a_n$ គេអាចកំណត់
\[
\begin{aligned}
a_2 &= 1 + 4a_1 = 1 + 4\\
a_3 &= 1 + 4a_2 = 1 + 4 \times (1 + 4) = 1 + 4 + 4^2\\
a_4 &= 1 + 4a_3 = 1 + 4 \times (1 + 4 + 4^2) = 1 + 4 + 4^2 + 4^3\\
&\ \,\vdots
\end{aligned}
\]
នោះ $a_n = 1 + 4 + 4^2 + 4^3 + \ldots + 4^{n-1}$ ។ គេបាន
\[ a_n = 1 + 4 + 4^2 + 4^3 + \ldots + 4^{n-1} = \frac{1 \times (4^n - 1)}{4 - 1} = \frac{4^n - 1}{3} \text{ ។} \]
ដូច្នេះ តួទី $n$ នៃស៊ីត $(a_n)$ គឺ $a_n = \dfrac{4^n - 1}{3}$ ។

\textbf{របៀបទី ២} ៖ ចំពោះ $n \ge 2$ គេបាន $a_{n+1} = 1 + 4a_n$, $a_n = 1 + 4a_{n-1}$ នោះ
\[
\begin{gathered}
\begin{aligned}
a_{n+1} - a_n &= 1 + 4a_n - (1 + 4a_{n-1})\\
&= 4(a_n - a_{n-1})
\end{aligned}\\
(1)
\end{gathered}
\]
បើគេតាង $b_n = a_{n+1} - a_n$ នោះពី $(1)$ គេបាន $b_n = 4b_{n-1} \Rightarrow \dfrac{b_n}{b_{n-1}} = 4$ នោះស៊ីត $(b_n)$ ជាស៊ីតធរណីមាត្រដែលមានផលធៀបរួមស្មើនឹង $4$ និងមានតួទី $1$ និងតួទី $n$ គឺ $b_1 = a_2 - a_1 = 1 + 4 \times 1 - 1 = 4$, $b_n = 4 \times 4^{n-1} = 4^n$ ។

ស៊ីត $(b_n)$ ជាផលសងស៊ីត $(a_n)$ តាមរូបមន្ត ចំពោះ $n \ge 2$ គេបាន
\[
\begin{aligned}
a_n &= a_1 + \sum_{k=1}^{n-1} 4^k = 1 + 4^1 + 4^2 + \ldots + 4^{n-1}\\
&= \frac{1 \times (4^n - 1)}{4 - 1} = \frac{4^n - 1}{3} \text{ ។}
\end{aligned}
\]
ដូច្នេះ តួទី $n$ នៃស៊ីត $(a_n)$ គឺ $a_n = \dfrac{4^n - 1}{3}$ ។

\textbf{របៀបទី ៣} ៖ គេប្រើស៊ីតជំនួយ $(r_n)$ រួចធ្វើផលដក $a_{n+1} - r_n$ ។

គេឱ្យ $r_n = -\dfrac{1}{3}$ គេបាន $a_{n+1} + \dfrac{1}{3} = 1 + 4a_n + \dfrac{1}{3} = 4a_n + \dfrac{4}{3} = 4 \cdot \left(a_n + \dfrac{1}{3}\right)$ ។

តាង $V_{n+1} = a_{n+1} + \dfrac{1}{3}$ ហើយ $V_n = a_n + \dfrac{1}{3}$ \quad $(1)$

គេបាន $V_{n+1} = 4V_n$ ហើយ $V_n = 4^{n-1} \cdot V_1$ ។ គេអាចគណនា $V_1$ តាមទំនាក់ទំនង $(1)$ គឺ $V_1 = a_1 + \dfrac{1}{3} = \dfrac{4}{3}$ ។

គេបាន $V_n = 4^{n-1} \cdot \dfrac{4}{3} = \dfrac{1}{3} \cdot 4^n$ ។ គេអាចគណនា $a_n$ តាមទំនាក់ទំនង $(1)$ គឺ $a_n = V_n - \dfrac{1}{3} = \dfrac{1}{3} \cdot 4^n - \dfrac{1}{3} = \dfrac{4^n - 1}{3}$ ។

\begin{note}
ដើម្បីកំណត់រកតម្លៃនៃស៊ីតជំនួយ $(r_n)$ គេតាង $r_n = \alpha \cdot n + \beta$ ($\alpha, \beta$ ថេរ) ជាស៊ីតដែលផ្ទៀងផ្ទាត់ $a_{n+1} = 1 + 4a_n$ ។ ជំនួសស៊ីតជំនួយទៅក្នុងស៊ីត $(a_n)$ គេបានសមីការដែលអាចឱ្យ គេកំណត់រក $\alpha$ និង $\beta$ ។
\[
\begin{aligned}
a_{n+1} &= 1 + 4a_n \iff \alpha \cdot (n + 1) + \beta\\
&= 1 + 4 \cdot (\alpha \cdot n + \beta) \text{ ឬ } 3 \cdot \alpha \cdot n + 3 \cdot \beta - \alpha + 1 = 0
\end{aligned}
\]
គេបាន $\begin{cases} 3 \cdot \alpha = 0\\ 3 \cdot \beta - \alpha + 1 = 0 \end{cases}$ ហើយ $\alpha = 0$, $\beta = -\dfrac{1}{3}$ ។

ដូច្នេះ $r_n = -\dfrac{1}{3}$ ដូច្នេះស៊ីត $(r_n)$ ហៅថា \textbf{ស៊ីតជំនួយ} ។
\end{note}

\begin{formulaBox}[ជាទូទៅ]
គេប្រើស៊ីតជំនួយ $(r_n)$ ដែល $r_n = \alpha \cdot n + \beta$ ($\alpha, \beta$ ថេរ) ដែលផ្ទៀងផ្ទាត់ទំនាក់ទំនងនៃស៊ីត ។ ក្នុងករណីនេះ គេជំនួស $(r_n)$ ក្នុងទំនាក់ទំនងនៃស៊ីត រួចគណនា តម្លៃ $\alpha$ និង $\beta$ បន្ទាប់មកកំណត់បានស៊ីត $(r_n)$ គេនឹងធ្វើផលសង $a_{n+1} - r_n$ ។
\end{formulaBox}

\begin{example}[លំហាត់គំរូ ២]
គេឱ្យស៊ីត $(a_n)$ កំណត់ដោយ $a_1 = 3$, $a_{n+1} = 2a_n - n + 1$ ។ កំណត់តួទី $n$ ។
\end{example}
\ans
តាង $(r_n)$ ដែល $r_n = \alpha n + \beta$ ជាស៊ីតជំនួយ ។ គេបាន
\[
\begin{aligned}
a_{n+1} = 2a_n - n + 1 &\iff \alpha \cdot (n + 1) + \beta\\
&= 2(\alpha \cdot n + \beta) - n + 1\\
&\iff n \cdot (\alpha - 1) + \beta - \alpha + 1 = 0
\end{aligned}
\]
គេបានប្រព័ន្ធ $\begin{cases} \alpha - 1 = 0\\ \beta - \alpha + 1 = 0 \end{cases}$ ហើយ $\alpha = 1$, $\beta = 0$ ។

នោះ $(r_n) = n$ ធ្វើផលសងរវាងស៊ីត $a_{n+1} - r_n$ គេបាន ៖
\[
\begin{gathered}
a_{n+1} - n = 2a_n - n + 1 - n = 2a_n - 2n + 1,\\
a_{n+1} - (n + 1) = 2(a_n - n)
\end{gathered}
\]
តាង $V_n = a_n - n$ \quad $(1)$ នោះ $V_{n+1} = a_{n+1} - (n + 1)$ គេបាន $V_{n+1} = 2 \cdot V_n$ ហើយ $V_n = 2^{n-1} \cdot V_1$ ។

គេគណនា $V_1$ តាម $V_n = a_n - n$, $V_1 = a_1 - 1 = 3 - 1 = 2$, $V_n = 2^{n-1} \cdot 2 = 2^n$ ។

គេអាចគណនា $V_n$ តាម $(1)$ គឺ $V_n = a_n - n$ នោះ $a_n = V_n + n = 2^n + n$ ។

ដូច្នេះ $a_n = 2^n + n$ ។

\begin{pro}
កំណត់តួទី $n$ នៃស៊ីត $(a_n)$ កំណត់ដោយ $a_1 = 5$, $a_n = 2a_{n-1} - n$ ។
\ans
$a_n = 2^n + (n + 2)$ ។
\end{pro}

\section{ទំនាក់ទំនងរវាង \texorpdfstring{$a_n$}{an} និង \texorpdfstring{$S_n$}{Sn}}

\begin{example}[ឧទាហរណ៍]
គេមានផលបូក $n$ តួដំបូងនៃស៊ីត $(a_n)$ គឺ $S_n = n^2 + n$ ។ កំណត់តួទី $n$ នៃស៊ីត $(a_n)$ ។

គេមានផលបូក $n$ តួដំបូងនៃស៊ីត $(a_n)$ គឺ
\[ S_n = n^2 + n \iff S_n = a_1 + a_2 + a_3 + \ldots + a_{n-1} + a_n = n^2 + n \]
\[ \Rightarrow S_{n-1} = a_1 + a_2 + a_3 + \ldots + a_{n-1} = (n - 1)^2 + (n - 1) \text{ ។} \]
ចំពោះ $n \ge 2$ គេបាន $S_n - S_{n-1} = a_n = n^2 + n - [(n - 1)^2 + (n - 1)]$
\[
\begin{gathered}
\begin{aligned}
a_n &= n^2 + n - (n - 1)^2 - n + 1\\
&= n^2 + n - n^2 + 2n - 1 - n + 1 = 2n
\end{aligned}\\
\text{នោះ } a_n = 2n \text{ ។}
\end{gathered}
\]
ចំពោះ $n = 1$ គេបាន $S_1 = 1^2 + 1 = 2$ ពិត ព្រោះ ($S_1 = a_1 = 2$) ។

ដូច្នេះ $a_n = 2n$ ។
\end{example}

ឧបមាថា $S_n$ គឺជាផលបូក $n$ តួដំបូងនៃស៊ីត $(a_n)$ នោះ $S_n = a_1 + a_2 + \ldots + a_{n-1} + a_n$ និង $S_{n-1} = a_1 + a_2 + \ldots + a_{n-1}$ ។

ចំពោះ $n \ge 2$ គេបាន $S_n - S_{n-1} = a_n$ ។ ម្យ៉ាងទៀត $S_1 = a_1$ ។

\begin{formulaBox}[ជាទូទៅ]
ចំពោះ $n \ge 2$ គេបាន $S_n - S_{n-1} = a_n$ និង $S_1 = a_1$ ។
\end{formulaBox}

\begin{note}
$S_n$ និង $S_{n+1}$ មានភាពចាំបាច់ដើម្បីរក $a_{n+1}$ គឺ $a_{n+1} = S_{n+1} - S_n$ ។
\end{note}

\begin{example}[លំហាត់គំរូ]
គេមានផលបូក $n$ តួដំបូងនៃស៊ីត $(a_n)$ គឺ $S_n = n^3 + 2n$ ។ កំណត់តួទី $n$ នៃស៊ីត $(a_n)$ ។
\end{example}
\ans
គេមានផលបូក $n$ តួដំបូងនៃស៊ីត $(a_n)$ គឺ $S_n = n^3 + 2n \Rightarrow S_{n-1} = (n - 1)^3 + 2(n - 1)$
\[
\begin{aligned}
a_n = S_n - S_{n-1} &= n^3 + 2n - (n - 1)^3 - 2(n - 1)\\
&= n^3 + 2n - 2n + 2 - n^3\\
&\quad + 3n^2 - 3n + 1\\
&= 3n^2 - 3n + 3 \text{ ។}
\end{aligned}
\]
ដូច្នេះ តួទី $n$ នៃស៊ីត $(a_n)$ គឺ $a_n = 3n^2 - 3n + 3$ ។

\begin{pro}
គេមាន $S_n = 4 - a_n - \dfrac{1}{2^{n-2}}$ $(n = 1, 2, 3, \ldots)$ ជាផលបូក $n$ តួដំបូងនៃស៊ីត $(a_n)$ ។
\begin{enumerate}[a]
  \item កំណត់ទំនាក់ទំនងកំណើនរវាង $a_{n+1}$ និង $a_n$ ។
  \item កំណត់តួទី $n$ នៃស៊ីត $(a_n)$ ។
\end{enumerate}
\ans
\begin{enumerate}[a]
  \item $2a_{n+1} = a_n + \dfrac{1}{2^{n-1}}$
  \item $a_n = \dfrac{2n}{2^n}$ ។
\end{enumerate}
\end{pro}

\section{ទំនាក់ទំនងកំណើនតួក្នុងទម្រង់ \texorpdfstring{$a_{n+2} + pa_{n+1} + qa_n = 0$}{a(n+2)+pa(n+1)+qa(n)=0}}

\begin{example}[ឧទាហរណ៍]
កំណត់តួទី $n$ នៃស៊ីត $(a_n)$ ដោយ $a_1 = 1$, $a_2 = 3$ ដែល $a_{n+2} = 3a_{n+1} - 2a_n$ $(n = 1, 2, \ldots)$ ។

គេស្គាល់ $a_1$, $a_2$ គេអាចគណនាតួ $a_3$ និងតួបន្តបន្ទាប់ដើម្បីពិនិត្យមើលលំនាំគំរូ
\[
\begin{aligned}
a_3 &= 3 \cdot a_2 - 2 \cdot a_1 = 9 - 2 = 7\\
a_4 &= 3 \cdot a_3 - 2 \cdot a_2 = 21 - 6 = 15\\
a_5 &= 3 \cdot a_4 - 2 \cdot a_3 = 45 - 14 = 31\\
a_6 &= 3 \cdot a_5 - 2 \cdot a_4 = 93 - 30 = 63
\end{aligned}
\]
គណនាផលសងតួនៃស៊ីតនេះ
\[
\begin{gathered}
a_2 - a_1 = 2, \quad a_3 - a_2 = 2^2,\\
a_4 - a_3 = 2^3, \quad \ldots,\\
a_n - a_{n-1} = 2^{n-1}
\end{gathered}
\]
\[
\begin{gathered}
a_n - a_1 = 2 + 2^2 + 2^3 + \ldots + 2^{n-1} \quad \text{នោះ}\\
a_n = 1 + 2 + 2^2 + 2^3 + \ldots + 2^{n-1}
\end{gathered}
\]
ជាផលបូកតួនៃស៊ីតធរណីមាត្រដែលមាន $1$ ជាតួទី $1$ និងផលធៀបរួមស្មើនឹង $2$ គេបាន $a_n = 2^n - 1$ ។

ជាទូទៅគេអាចគណនា $a_n$ តាមវិធីខាងក្រោម ៖
\[
\begin{aligned}
&a_{n+2} = 3a_{n+1} - 2a_n\\
&\iff a_{n+2} - 2a_{n+1} = a_{n+1} - 2a_n\\
&= a_n - 2a_{n-1} = a_{n-1} - 2a_{n-2} = \ldots\\
&= a_2 - 2a_1 = 3 - 2 \times 1 = 1
\end{aligned}
\]
នោះ $a_{n+2} - 2a_{n+1} = a_{n+1} - 2a_n = 1$ \quad $(1)$ ។

ម្យ៉ាងទៀត $a_{n+2} = 3a_{n+1} - 2a_n \iff a_{n+2} - a_{n+1} = 2(a_{n+1} - a_n)$ ។

តាង $b_n = a_{n+1} - a_n$ នោះ $(b_n)$ ឬ $(a_{n+1} - a_n)$ ជាស៊ីតធរណីមាត្រដែលមានផលធៀបរួមស្មើនឹង $2$ និងតួទី $1$ ស្មើនឹង $b_1 = a_2 - a_1 = 3 - 1 = 2$ ហើយ $b_n = 2 \cdot (2)^{n-1} = 2^n$ \quad $(2)$

តាម $(1)$ និង $(2)$ គេបាន
\[
\begin{cases}
a_{n+1} - 2a_n = 1 & (1)\\
a_{n+1} - a_n = 2^n & (2)
\end{cases}
\]
ដក $(1) - (2)$ គេបាន $-a_n = 1 - 2^n$ ។

ដូច្នេះ តួទី $n$ នៃស៊ីត $(a_n)$ គឺ $a_n = 2^n - 1$ ។
\end{example}

គេបានរបៀបរកតួទូទៅតាមទំនាក់ទំនងដូចខាងក្រោម ៖

\begin{formulaBox}[ជាទូទៅ]
\textbf{របៀបដោះស្រាយ $a_{n+2} + pa_{n+1} + qa_n = 0$}

បើទំនាក់ទំនងនៃស៊ីត $(a_n)$ មានទម្រង់ $a_{n+2} + pa_{n+1} + qa_n = 0$ គេពិចារណាពីរចំនួន $\alpha$ និង $\beta$ ដែលបំពេញលក្ខខណ្ឌ $\alpha + \beta = -p$, $\alpha\beta = q$ នោះគេបាន $a_{n+2} - (\alpha + \beta)a_{n+1} + \alpha\beta a_n = 0$ ។

ទំនាក់ទំនងនេះអាចបម្លែងតាមរបៀបពីរយ៉ាងដូចខាងក្រោម ៖
\[
\begin{gathered}
a_{n+2} - \alpha a_{n+1} = \beta(a_{n+1} - \alpha a_n)\\
a_{n+2} - \beta a_{n+1} = \alpha(a_{n+1} - \beta a_n)
\end{gathered}
\]
ទំនាក់ទំនងទាំង $2$ ខាងលើនេះ គឺជាស៊ីតធរណីមាត្រអាចឱ្យគេរក $a_n$ បាន ។

គេអាចរក $\alpha$ និង $\beta$ ដោយដោះស្រាយសមីការ $x^2 + px + q = 0$ ព្រោះគេស្គាល់ផលបូកនិងផលគុណវា ។

សមីការ $x^2 + px + q = 0$ ហៅថា ``\textbf{សមីការសម្គាល់}នៃ $a_{n+2} - (\alpha + \beta)a_{n+1} + \alpha\beta a_n = 0$'' ។
\end{formulaBox}

\begin{example}[លំហាត់គំរូ]
ស៊ីត $(a_n)$ កំណត់ដោយ $a_1 = 1$, $a_2 = 13$, $a_{n+2} = a_{n+1} + 6a_n$ $(n = 1, 2, \ldots)$ ។ កំណត់តួទី $n$ នៃស៊ីត ។
\end{example}
\ans
កំណត់តួទី $n$ នៃស៊ីត $(a_n)$ ៖ $a_{n+2} = a_{n+1} + 6a_n \iff a_{n+2} - a_{n+1} - 6a_n = 0$

សមីការសម្គាល់ ៖ $x^2 - x - 6 = 0$ សមីការមានឫស $x = 3$ ឬ $x = -2$ ។ គេបាន
\[
\begin{cases}
a_{n+2} + 2a_{n+1} = 3(a_{n+1} + 2a_n) & (i)\\
a_{n+2} - 3a_{n+1} = -2(a_{n+1} - 3a_n) & (ii)
\end{cases}
\]
$(i)$ \quad $a_{n+2} + 2a_{n+1} = 3(a_{n+1} + 2a_n)$ ។ តាង $b_n = a_{n+1} + 2a_n$ នោះស៊ីត $(b_n)$ ជាស៊ីតធរណីមាត្រដែលមានផលធៀបរួមស្មើនឹង $3$ ហើយតួទី $1$ គឺ $b_1 = a_2 + 2a_1 = 13 + 2 \times 1 = 15$ និង $b_n = 15 \cdot (3)^{n-1}$ ។ គេបាន $a_{n+1} + 2a_n = 15 \cdot (3)^{n-1}$ \quad $(1)$ ។

$(ii)$ \quad $a_{n+2} - 3a_{n+1} = -2(a_{n+1} - 3a_n)$ ។ តាង $c_n = a_{n+1} - 3a_n$ នោះស៊ីត $(c_n)$ ជាស៊ីតធរណីមាត្រដែលមានផលធៀបរួមស្មើនឹង $-2$ ហើយតួទី $1$ គឺ $c_1 = a_2 - 3a_1 = 13 - 3 \times 1 = 10$ និង $c_n = 10 \cdot (-2)^{n-1}$ ។ គេបាន $a_{n+1} - 3a_n = 10 \cdot (-2)^{n-1}$ \quad $(2)$ ។

តាម $(1)$ និង $(2)$ គេបាន
\[
\begin{cases}
a_{n+1} + 2a_n = 15 \cdot (3)^{n-1} & (1)\\
a_{n+1} - 3a_n = 10 \cdot (-2)^{n-1} & (2)
\end{cases}
\]
ដក $(1) - (2)$ ៖ $5a_n = 15(3)^{n-1} - 10(-2)^{n-1}$
\[ a_n = 3 \times 3^{n-1} - 2 \times (-2)^{n-1} = 3^n + (-2)^n \]
ដូច្នេះ តួទី $n$ នៃស៊ីត $(a_n)$ គឺ $a_n = 3^n + (-2)^n$ ។

\begin{pro}
ស៊ីត $(a_n)$ កំណត់ដោយទំនាក់ទំនង $a_1 = 0$, $a_2 = 3$, $a_{n+2} = 4a_{n+1} - 3a_n$ $(n = 1, 2, \ldots)$ ។ កំណត់តួទី $n$ នៃស៊ីត $(a_n)$ ។
\ans
$a_n = \dfrac{1}{2}(3^n - 3)$ ។
\end{pro}

\section*{លំហាត់អនុវត្តន៍}
\begin{enumerate}[1]
  \item ស៊ីត $(a_n)$ កំណត់ដោយទំនាក់ទំនងកំណើនដូចខាងក្រោម ៖
  \begin{enumerate}[a]
    \item $a_1 = 3$, $a_{n+1} = 2a_n - 4$
    \item $a_1 = 5$, $a_{n+1} = 3a_n - 2n$ ។
  \end{enumerate}
  កំណត់តួទី $4$ នៃស៊ីតនេះ ។
  \item ស៊ីត $(a_n)$ កំណត់ដោយទំនាក់ទំនងកំណើន $a_1 = 1$, $a_2 = 2$, $a_{n+2} = a_{n+1} + a_n$ ។ កំណត់តួទី $7$ នៃស៊ីតនេះ ។
  \item ស៊ីត $(a_n)$ កំណត់ដោយទំនាក់ទំនងកំណើន $a_1 = 1$, $a_{n+1} = pa_n + q$ ។ គណនាតម្លៃ $p$ និង $q$ បើគេដឹងតួទី $3$ ស្មើនឹង $6$ និងតួទី $5$ ស្មើនឹង $86$ ។
  \item កំណត់តួទូទៅនៃស៊ីត $(a_n)$ កំណត់ដោយទំនាក់ទំនងកំណើនដូចខាងក្រោម ៖
  \begin{enumerate}[a]
    \item $a_1 = 1$, $3a_{n+1} = 2a_n + 3$ \quad $(n = 1, 2, \ldots)$
    \item $a_1 = 1$, $3a_{n+1} = 3a_n + 4$ \quad $(n = 1, 2, \ldots)$
    \item $a_1 = 1$, $a_{n+1} = a_n + n$ \quad $(n = 1, 2, \ldots)$ ។
  \end{enumerate}
  \item គេមានស៊ីត $(a_n)$ កំណត់ដោយ $a_1 = 4$, $a_{n+1} = \dfrac{4a_n - 9}{a_n - 2}$ $(n = 1, 2, 3, \ldots)$ ។
  \begin{enumerate}[a]
    \item ស្រាយបញ្ជាក់ថា $a_n \ne 3$ ចំពោះគ្រប់ $n$ ។
    \item យក $b_n = \dfrac{1}{a_n - 3}$ និងកំណត់តួទូទៅនៃស៊ីត $(b_n)$ ។ កំណត់តួទូទៅនៃស៊ីត $(a_n)$ ។
  \end{enumerate}
  \item គេមាន $S_n$ ជាផលបូក $n$ តួដំបូងនៃស៊ីត $(a_n)$ ។ បើ $S_n$ បំពេញលក្ខខណ្ឌ $S_n = 4 - a_n - \dfrac{1}{2^{n-2}}$ $(n = 1, 2, 3, \ldots)$ នោះចូរ ៖
  \begin{enumerate}[a]
    \item កំណត់ដោយទំនាក់ទំនងកំណើនរវាង $a_{n+1}$ និង $a_n$ ។
    \item កំណត់តួទី $n$ នៃស៊ីត $(a_n)$ ។
  \end{enumerate}
  \item គេមាន $S_n$ ជាផលបូក $n$ តួដំបូងនៃស៊ីត $(a_n)$ ហើយបំពេញលក្ខខណ្ឌ $(a_n)$ ៖ $a_1 = 1$, $S_n = a_{n+1} + n^2$ $(n \ge 1)$ ។ កំណត់តួទី $n$ នៃស៊ីត $(a_n)$ ។
  \item គេមាន $S_n$ ជាផលបូក $n$ តួដំបូងនៃស៊ីត $(a_n)$ ហើយ $S_n$ បំពេញលក្ខខណ្ឌ $S_n = \dfrac{n}{n-1} \cdot a_n$, $n \ge 2$ ។
  \begin{enumerate}[a]
    \item បញ្ជាក់រក $a_n$ $(n \ge 3)$ អនុគមន៍នឹង $n$ និង $a_{n-1}$ ។
    \item បញ្ជាក់រក $S_n$ $(n \ge 2)$ អនុគមន៍នឹង $n$ និង $S_{n-1}$ ។
    \item ឧបមាថា $a_1 = 1$ រកតួទី $n$ នៃស៊ីត $S_n$ ដែល $n \ge 1$ ។
  \end{enumerate}
  \item គេមានស៊ីត $(a_n)$ កំណត់ដោយទំនាក់ទំនងកំណើន $(a_n)$ ៖ $a_1 = 2$, $a_{n+1} = \dfrac{a_n}{a_n + 3}$ $(n = 1, 2, 3, \ldots)$ ។
  \begin{enumerate}[a]
    \item តាង $b_n = \dfrac{1}{a_n}$ ។ កំណត់ទំនាក់ទំនងកំណើនរវាង $b_n$ និង $b_{n+1}$ ។
    \item កំណត់តួទី $n$ នៃស៊ីត $a_n$ ។
  \end{enumerate}
\end{enumerate}
